En y repensant, il y a quelque chose d'à la fois amusant et cohérent à finaliser cette thèse. J'ai débuté cette aventure, convaincu que c'était la voie professionnelle qu'il me fallait alors que, trois ans auparavant, je jurais par tous les saints que je ne ferais jamais de thèse. Puis, je repense au gamin que j'étais en primaire, écrivant fièrement sur une copie de rédaction qu'il serait docteur (en médecine) et chercheur scientifique. Alors, me voici aujourd'hui docteur, mais clairement pas en médecine\footnote{Désolé, Maman !}. Bien que je ne sache pas encore si je serai chercheur toute ma vie, je me dis, après ces trois belles années à faire de la recherche, que cela vaut bien le coup de continuer encore un peu.

Ces trois années et demie (soyons précis) ont été marquées par des moments de doute, un syndrome de l'imposteur, des déceptions et des expériences infructueuses..., mais elles ont aussi abouti à ce tapuscrit, à quelques publications et, surtout, à de belles rencontres et de beaux souvenirs. Il me paraît donc tout à fait naturel et juste de commencer ce tapuscrit par des remerciements sincères à toutes ces personnes qui m'ont soutenu et accompagné dans cette aventure.

% Remerciements encadrants thèse
Lorsqu'on me demande sur quoi repose le bon déroulement d'une thèse, je réponds toujours : \textit{"De bons encadrants !"}. Et je pense pouvoir affirmer avoir eu de bons encadrants, car cette thèse et toutes ces contributions n'auraient jamais vu le jour sans eux. Par ces quelques mots, je tiens à leur exprimer toute ma gratitude pour leur encadrement, leur soutien, et l'autonomie qu'ils m'ont accordée tout au long de cette aventure.  
Merci à Bertrand Mathieu de m'avoir fait confiance dès le début pour ce sujet. Ton accompagnement, autant humain que professionnel, m'a permis de grandir dans ce parcours exigeant. J'ai énormément appris grâce à tes remarques franches et constructives, ta pédagogie et ta bienveillance, qui m'ont guidé dans mon apprentissage du métier de chercheur. Merci à Abdelkader Lahmadi pour ton soutien, toujours empreint de bienveillance, et pour tes observations pertinentes qui ont énormément enrichi mes travaux. Un immense merci également à Raouf Boutaba et Isabelle Chrisment. Malgré vos agendas bien chargés, vous avez toujours su être disponibles au moment où c'était nécessaire pour faire avancer cette thèse.

% Remerciements jury de thèse  
Je souhaite également remercier chaleureusement Sandrine Vaton et Nadjib Aitsaadi pour avoir accepté d’être les rapporteurs de cette thèse et pour le temps consacré à l’évaluation de mon travail. Merci à Hind Taleb, Yacine Ghamri et Claudia Ignat pour leur participation au jury. Merci à tous pour votre intérêt pour mes travaux. Merci aussi à Ye-Qiong Song, qui avec Claudia, ont suivi mes travaux depuis le début et m’ont apporté leur accompagnement et leurs précieux conseils.

% Remerciements S. Tuffin  
Un merci tout particulier à Stéphane Tuffin, dont l'engagement a été essentiel pour cette thèse. Merci pour ta franchise, tes conseils avisés et ton soutien logistique. Tu n'as ménagé aucun effort pour m'aider à obtenir les ressources nécessaires à mes expériences, et ton implication fut d'une très grande aide.

% Remerciements MOSAICO  
J’ai également eu la chance de réaliser cette thèse au sein du projet ANR MOSAICO, où j’ai rencontré des doctorants, ingénieurs, et chercheurs très compétents. Merci à Philippe, Hichem, Marius, Xavier, Thibault, Guillaume, Huu Nghia, Edgardo : les discussions, les échanges, les réunions, les Orange Open Tech Days 2023, et même les moments les plus informels comme les restaurants ou les \textit{social events}, ont rendu ce parcours doctoral aussi enrichissant scientifiquement que socialement. Grâce à cette émulation collective, j’ai pu plonger rapidement dans mon sujet et avancer plus facilement dans mes travaux.

% Remerciements Orange: Olivier, Delphine, + team ITEQ  
Je souhaite remercier mes collègues de l'équipe ITEQ à Orange Innovation Lannion, merci pour votre accueil chaleureux. Je me suis senti intégré dans l'équipe dès les premiers jours et les (longues) pauses café (et surtout les \textit{PPVR}) de l'équipe vont me manquer. Un grand merci à Olivier Carli, en particulier, pour son implication dans ma thèse, depuis le recrutement, mon arrivée dans l'équipe jusqu'à l'organisation de mes déplacements en conférence ou en formation. Ton aide fut précieuse.

% Remerciements co-auteurs Orange: Minqi, Nicolas, Eric  
Merci aussi à Minqi, Eric et Nicolas pour votre aide lors des tests au labo FA. Ce fut assez laborieux à mettre en place, entre les nouveaux bugs qui apparaissent dès qu'on a résolu les anciens et les allers-retours sur Pégase. Merci de votre aide, car une partie de ces travaux n'aurait pas été réalisée sans vous.

Je tiens également à remercier chaleureusement les chercheur·e·s, ingénieur·e·s et tout le personnel d’Orange Innovation Lannion, qui ont contribué, directement ou indirectement, au succès de cette thèse\footnote{Merci surtout à la cantine Sodexo pour vos plats qui m'épargnaient la charge mentale de la cuisine.}. J’ai toujours trouvé des personnes disponibles et prêtes à partager leur expertise sur de nombreux sujets tels que les réseaux mobiles ou Wi-Fi, la science des données ou encore pour m'aider à me dépatouiller avec le proxy d'Orange. J'ai pris plaisir à réaliser ma thèse sur le site de Lannion, et merci à chacun d’avoir pris le temps de m’écouter, de répondre à mes questions, et également de votre soutien lors de ma participation au concours \textit{Ma thèse en 3 minutes}.   


Merci au LORIA et en particulier aux membres de l'équipe RESIST pour leur accueil lors des derniers mois de ma thèse pour finaliser la rédaction du tapuscrit.

% Remerciements doctorants Orange: Fatiha, Ziad, Abraham, Diane, Thomas, Lenaïg, Vercine  
Je ne saurais conclure sans adresser un immense merci à celles et ceux avec qui j'ai partagé mon quotidien (et, avouons-le surtout les galères) durant une bonne partie de ces trois années de thèse. Merci à Ziad mon voisin de bureau pour ta gentillesse (et les pâtisseries libanaises que tu rapportais), à Fatiha pour ta compagnie et nos discussions sur One Piece durant la première année (et surtout les \textit{msemmen}), à Laure pour ton aide après mon opération, à Diane pour les soirées organisées (et ton incroyable absence de second degré), à Thomas qui a toujours la réponse juste à tout problème (et aussi la blague douteuse/contrepèterie appropriée), à Hassina pour ta gentillesse (et aussi tes tiramisus), à Lenaïg pour tes râleries sur les menus végés de la cantine (et tes tacles souvent gratuits), à Vercine pour tes réactions toujours hilarantes et à Régis le geek Lannionnais aux mille et une refs. Merci aussi à Abraham, Marco, Othmane, Khaled, Rim et aux doctorant·e·s et stagiaires rencontrés à Lannion. Merci à tous pour les rires, les potins, les confidences, les dramas, les plaintes sur le doctorat, les parties interminables de Wizard, de ping-pong, les soirées jeux de société, les afterworks, les randos et même ce que je n'ai pas réussi à oublier\footnote{Je parle de ces parties de baby-foot assaisonnées des piques de la bretonne.}, qui ont apporté un peu de soleil et aidé à apprécier le \textit{si beau temps} de Lannion.

%Merci aussi à mes ami·e·s et camarades qui m'ont apporté leur soutien sous des formes diverses durant cette thèse.

% Remerciements famille  
Enfin, je tiens à remercier du fond du cœur ma famille, mon père, ma mère et mes frères. Il aurait été plus logique de commencer ces remerciements par vous, tant votre soutien et vos encouragements ont été importants pour moi. Merci pour vos conseils, vous m'avez toujours poussé à faire de mon mieux, même si vous ne compreniez pas toujours ce que je faisais comme travaux de recherche, ni pourquoi je faisais un doctorat\footnote{J'avoue m'être moi aussi parfois posé la question.}. Promis, c'est mon dernier diplôme, après je commence à travailler\footnote{Sauf si je me décide à faire un CAP boulanger ou pâtissier, vu mon amour à peine dissimulé pour le pain et les pâtisseries.}.

% Remerciements Nouchan  
Je termine par toi, Axelle. Merci pour tes remarques toujours pertinentes, de m'avoir soutenu, encouragé et supporté pendant que je me plaignais de mon code, de la rédaction, des commentaires des reviewers ou juste comme ça. Ta patience, ta compréhension et ton amour me rendent tout beaucoup plus facile et juste pour cela merci d'être toujours à mes côtés.
